\documentclass[11pt,a4paper]{article}
\usepackage[utf8]{inputenc}
\usepackage[T1]{fontenc}
\usepackage{lmodern}
\usepackage[margin=1in]{geometry}
\usepackage[hyphens]{url}
\usepackage{microtype}
\usepackage{parskip}

\begin{document}
\thispagestyle{empty}

\today

\vspace{0.8em}

\noindent The Editor-in-Chief\\
\textit{Journal of Imaging}, MDPI

\vspace{0.8em}

\noindent Dear Editor,

\vspace{0.5em}

We are pleased to submit our manuscript, \textbf{``HypEReg: Hyperelastic Regularization for Reducing Folding in Transformer-Based Brain MRI Registration''}, for consideration as an original research article in the \textit{Journal of Imaging}.

\vspace{0.5em}
\noindent\textbf{Why this work matters to the medical image registration community.}
Transformer-based registration has advanced overlap accuracy, but many models still output deformation fields with non-positive Jacobian regions, limiting reliability for morphometry and atlas-based downstream pipelines. Our work addresses this core bottleneck by introducing \textbf{HypEReg}, a determinant-level hyperelastic regularizer that can be added to dense-flow Transformer training as a loss-side plug-in (no architectural change at inference).

\vspace{0.5em}
\noindent\textbf{Main technical contributions.}
\begin{itemize}
\item We propose a compact Jacobian-aware objective that combines a non-linear volume-distortion term and an explicit anti-folding hinge, directly supervising \(\det J_\phi\).
\item We show that this regularization strategy is practical for community adoption: it is plug-and-play for Transformer dense-flow backbones and introduces \emph{zero inference overhead}.
\item We provide full subject-level inferential analysis (paired Wilcoxon + BH-FDR), bootstrap confidence intervals, and open reproducibility artifacts.
\end{itemize}

\vspace{0.2em}
\noindent\textbf{Key verified results (IXI, 115 held-out subjects).}
Compared with plain TransMorph, HypEReg-TransMorph improves grouped Dice (0.7537 vs.\ 0.7527), lowers HD95 (5.3234 vs.\ 5.6872\,mm), and lowers ASSD (1.3570 vs.\ 1.4073\,mm), while dramatically improving deformation validity: the non-positive Jacobian ratio drops from \(1.502\times10^{-2}\) to \(1.5\times10^{-5}\) (approximately three orders of magnitude; \(>1000\times\) reduction), and SDlogJ drops from 0.5064 to 0.3280. Runtime and memory remain identical to TransMorph (0.0822\,s/case, 5.685\,GB).

\vspace{0.2em}
\noindent\textbf{Downstream relevance and practical impact.}
Beyond registration-level metrics, we evaluate a downstream OASIS zero-shot multi-atlas fusion task (\(N_{\mathrm{atlas}}=6\), \(n=20\) targets). HypEReg-TransMorph attains the highest fused Dice (0.8271), outperforming TransMorph (0.8201), TransMorphBayes (0.8058), and MIDIR (0.7696). This demonstrates that large-scale reduction of fold-prone deformation behavior translates into more consistent label propagation and better fusion quality in a clinically relevant downstream setting.

\vspace{0.5em}
\noindent\textbf{Fit for \textit{Journal of Imaging}.}
We believe the manuscript aligns strongly with the journal's aims: it contributes a rigorous, reproducible methodological advance with immediate value for medical image analysis workflows that depend on physically plausible deformations. Code and reproducibility assets are openly available at \url{https://github.com/Zzxmh/HypEReg-TransMorph} and archived at Zenodo (DOI: 10.5281/zenodo.19888526).

\vspace{0.5em}
\noindent\textbf{Submission declarations.}
This manuscript is original, has not been published, and is not under consideration elsewhere. All authors have approved the submission. The authors declare no conflict of interest, and no external funding was received.

\vspace{1em}

\noindent Sincerely,\\[0.5em]
\noindent Shiyi Xu (on behalf of all authors)\\
Capital Medical University, Beijing 100069, China\\
\texttt{xsy4831@mail.ccmu.edu.cn}

\vspace{0.5em}
\noindent Erjin Zhou (corresponding author)\\
Dexmal, Beijing, China\\
\texttt{zhouerjin@gmail.com}

\end{document}
